% Responsable(s): por asignar
% Objetivo: explicar jerarquía de memoria y definir latencia, tiempo de acceso,
% ancho de banda y otras métricas necesarias para interpretar los resultados.

% Redacción pendiente.
\subsection{Jerarquía de memoria}

La jerarquía de memoria busca reducir la diferencia de velocidad existente
entre el procesador y la memoria principal mediante distintos niveles de 
almacenamiento. Los niveles más cercanos al procesador proveen tiempos 
de acceso menores, mientras que la memoria principal ofrece una mayor capacidad
a costa de una latencia superior. Por esta razón, los sistemas modernos 
usan la memoria caché para asi mantener cerca del procesador los datos \cite{son2013reducing}.

Sin embargo, cuando los datos solicitados no se encuentran en los niveles
superiores, es necesario acceder a niveles inferiores y, posteriormente, 
a la memoria principal. Estos accesos pueden hacer periodos de espera en 
la ejecución del procesador. Park et al. muestran que los diferentes niveles 
de la jerarquía interactúan directamente con la CPU y que el nivel en el 
cual se produce un retraso puede afectar el rendimiento global del sistema 
\cite{park2022dvfs}.

\subsection{Latencia y tiempo de acceso}

La latencia de memoria representa el tiempo requerido para que una solicitud 
de acceso produzca los datos necesarios para asi continuar la ejecución.
En este trabajo, el tiempo de acceso se analiza a partir de esta latencia,
debido a que se debe determinar la rapidez con la que la memoria responde
a las solicitudes.

La latencia de una memoria DRAM depende de diferentes operaciones internas.
Entre ellas se encuentran la activación de filas, la precarga y los ciclos
periódicos de refresco.Estos mecanismos pueden impedir temporalmente nuevos 
accesos y aumentar el tiempo necesario para atender las solicitudes de memoria 
\cite{shahbazi2026ddr4}.

La latencia observada por el procesador tampoco depende únicamente del dispositivo
DRAM. Esmaili-Dokht et al. señalan que la latencia \textit{load-to-use} incluye 
también el tiempo que una solicitud permanece dentro del procesador, incluyendo 
elementos como la jerarquía de caché y la interconexión interna del chip. Por eso,
dos sistema que utilizan tecnología de memoria similar pueden presentar tiempos de 
acceso distintos \cite{esmailidokht2024mess}.

\subsection{Ancho de banda}

El ancho de banda de memoria representa la cantidad de datos que puede
transferirse durante un determinado intervalo de tiempo. Esta métrica permite
evaluar la capacidad del sistema de memoria para atender grandes volúmenes de
información.

Debe distinguirse entre el ancho de banda teórico de una tecnología y el
ancho de banda que puede alcanzarse en una implementación real. Operaciones
como los ciclos de refresco, la activación y la precarga pueden reducir el
throughput efectivo de una memoria DRAM. Con una implementación de cuatro memorias
DDR4-2400, Shahbazi y Tsui muestran cómo estos retardos limitan el aprovechamiento 
del ancho de banda \cite{shahbazi2026ddr4}. 

\subsection{Relación entre latencia y ancho de banda}

La latencia y el ancho de banda no son métricas independientes. Cuando la
cantidad de solicitudes de memoria aumenta, también aumenta la utilización
del ancho de banda disponible. Mientras el sistema se encuentra lejos de su
punto de saturación, la latencia puede mantenerse relativamente estable; sin
embargo, al aproximarse a la saturación, nuevas solicitudes generan mayor
contención y pueden producir incrementos significativos en el tiempo de
acceso \cite{esmailidokht2024mess}.

Los resultados obtenidos mediante el benchmark Mess muestran además que la
composición del tráfico influye sobre este comportamiento. El tráfico
compuesto únicamente por lecturas alcanza, en general, menor latencia y mayor
ancho de banda que configuraciones que incluyen escrituras, debido a las
restricciones temporales adicionales asociadas a estas operaciones
\cite{esmailidokht2024mess}.

\subsection{Memory wall y memory stalls}

La diferencia entre la velocidad de procesamiento y el tiempo necesario para
acceder a la memoria constituye uno de los problemas fundamentales en el
rendimiento de los sistemas computacionales. Aunque la capacidad y el ancho
de banda de DRAM han mejorado considerablemente, su latencia de acceso
aleatorio ha evolucionado de forma mucho más lenta, manteniendo la denominada
\textit{memory wall} \cite{son2013reducing}.

Una consecuencia de este problema son los \textit{memory stalls}, intervalos
durante los cuales la CPU debe esperar la resolución de una operación de
memoria antes de continuar procesando instrucciones. Estos periodos de espera
afectan métricas como las instrucciones por ciclo (IPC). Park et al. reportan
que una reducción de los ciclos de espera de memoria mediante una gestión
coordinada de la CPU y la jerarquía de memoria produjo una reducción del
26\% en los \textit{memory stalls per cycle} y un incremento del 11\% en IPC
en su plataforma experimental \cite{park2022dvfs}.
